%=============================================================================
\section{Problem Formulation}
\label{sec:problem}
%=============================================================================

\paragraph{Typed events.}
Let $\mathcal{C}=\{\text{Ethereum},\text{BSC},\text{Arbitrum}\}$.
A decoded action is a tuple $a=(t,c,u,v,y,\kappa,q,\pi)$ with timestamp $t$, chain $c\in\mathcal{C}$, source and destination addresses $u,v$, semantic type $y\in\mathcal{Y}$, token $\kappa$, log-scaled amount $q=\log(1+\text{amount})$ and protocol label $\pi$ (empty for plain transfers).
The type vocabulary is
\begin{align*}
\mathcal{Y}=\{&\textsf{transfer},\textsf{swap},\textsf{bridge\_dep},\textsf{bridge\_wdr},\textsf{lend\_dep},\\
&\textsf{lend\_wdr},\textsf{split},\textsf{merge},\textsf{mixer/exit},\textsf{other}\}.
\end{align*}

\paragraph{Trajectories and prefixes.}
Starting from a seed address $s$ named in a public incident report (or, for a negative, an address matched to the incident, Sect.~\ref{sec:data}), a trajectory $\tau=(a_1,\dots,a_n)$ is the time-ordered sequence of actions reached by bounded value-flow expansion (Sect.~\ref{sec:traj}).
The prefix of length $k$ is $\tau_{1:k}=(a_1,\dots,a_k)$; at wall-clock time $t$ the monitor has observed $\tau_{1:k(t)}$ with $k(t)=|\{i: t_i\le t\}|$.

\paragraph{Endpoint and lead time.}
For an incident trajectory whose exit is confirmed by a public source, the endpoint index $e^\star$ is the first action whose type is \textsf{bridge\_dep}, \textsf{lend\_dep} or \textsf{mixer/exit}, i.e.\ the first step after which the funds leave the monitored chain or enter a pool that breaks the trace; $t_e=t_{e^\star}$.
A detector outputs $\hat p_k=\Pr(\text{laundering}\mid\tau_{1:k})$ from $a_1,\dots,a_k$ only.
Given a threshold $\theta$, the first alert is at $k_a=\min\{k\ge2:\hat p_k\ge\theta\}$ and
\begin{equation}
\text{lead time}=t_e-t_{k_a},\qquad \text{lead steps}=e^\star-k_a .
\label{eq:lead}
\end{equation}
An alert with $k_a\ge e^\star$ is \emph{late}; late alerts are counted separately and never enter lead-time statistics.
An alert on a negative trajectory at any $k$ is a false alert, so false-alert rates are per trajectory, not per prefix.

\paragraph{Two views of ``early''.}
We evaluate two families of prefixes.
\emph{Online horizons} $k(t_1+h)$ for $h\in\{1\,\text{min},\dots,24\,\text{h}\}$ are available to a live monitor.
\emph{Relative checkpoints} $k=\lceil r n\rceil$, $r\in\{25,50,75,100\}\%$, require the final length $n$ and are therefore a retrospective diagnostic only; they are never used as a model input (a feature equal to $k/n$ was removed after we found it leaked $n$, Sect.~\ref{sec:features}).

\paragraph{Threat model and scope.}
The adversary controls pseudonymous addresses holding stolen or illicit funds and may use any public bridge, DEX, lending market or mixer.
The defender knows one seed address from a public report and observes public chain data only.
Breaking mixer anonymity, discovering incidents by whole-chain scanning, and attributing real-world identities are out of scope.

%=============================================================================
\section{The CAGI-ED Architecture}
\label{sec:method}
%=============================================================================

Figure~\ref{fig:arch} shows the four layers of CAGI-ED: (i) a multi-chain collector with an evidence cache, (ii) a semantic decoder that turns raw transfers into typed actions, (iii) a trajectory builder that follows tainted value from the seed, and (iv) an online scorer that recomputes prefix features on every new action, calibrates the score and emits an explained alert.

\begin{figure}[t]
\centering
\IfFileExists{figures/fig1_architecture.pdf}{\includegraphics[width=\linewidth]{figures/fig1_architecture.pdf}}{\fbox{\parbox[c][4.2cm][c]{0.95\linewidth}{\centering Figure 1 (architecture, drawn in draw.io): see \texttt{paper/figures/fig1\_architecture.pdf}}}}
\caption{CAGI-ED architecture. Raw explorer responses are cached with checksums, decoded into typed actions using a verified protocol map, and expanded from the seed into a value-flow trajectory. On every new action the scorer recomputes the prefix feature vector $\mathbf{x}(\tau_{1:k})$, applies the gradient-boosted model $f$ and Platt calibration, and raises an alert with its top TreeSHAP contributions and AADAPT technique IDs once $\hat p_k\ge\theta$.}
\label{fig:arch}
\end{figure}

\subsection{Collection and Evidence Cache}
For Ethereum and Arbitrum the collector calls the Etherscan V2 account endpoints (normal, internal and ERC-20 transactions) with block-range pagination; for BSC it calls NodeReal's \texttt{nr\_getAssetTransfers} in both directions.
Every response is stored verbatim with its request time and SHA-256 checksum under \emph{(chain, address, incident)}, so two incidents that touch the same infrastructure never overwrite each other's evidence and every feature value can be traced back to the bytes it was computed from.
The collection window is $[b_0, b_0+72\,\text{h}]$ from the incident's start block $b_0$.

\subsection{Semantic Decoding}
A protocol map lists contracts whose role was verified on-chain (by event signature or function selector) or in explorer labels: 9 bridge endpoints of 8 families, 6 DEX routers and pools, 2 mixers and 1 lending market.
Each raw leg becomes a canonical action; legs of the same transaction are then merged into one semantic action when the structure is unambiguous: two or more tokens moving in and out of the same address form a \textsf{swap}; one source paying two or more destinations forms a \textsf{split}; two or more sources paying one destination form a \textsf{merge}.
A deposit to a verified bridge (lending market, mixer) becomes \textsf{bridge\_dep} (\textsf{lend\_dep}, \textsf{mixer/exit}) regardless of the merge rule.
Duplicate legs returned by overlapping pages are removed on (chain, transaction hash, log index).

\subsection{Bounded Value-Flow Trajectories}
\label{sec:traj}
Let $A$ be the decoded actions of the processed addresses inside the window, sorted by time, and let $\rho(a)=\mathrm{e}^{q}-1$ recover the raw amount.
For a source $u$ and token $\kappa$ write $O(u,\kappa)=\sum_{a\in A:\,u_a=u,\kappa_a=\kappa}\rho(a)$.
The builder keeps a depth map $d(\cdot)$ with $d(s)=0$ and accepts action $a$ iff
\begin{equation}
d(u_a)<D,\qquad t_a-t_1\le H,\qquad \frac{\rho(a)}{O(u_a,\kappa_a)}\ge\eta \;\;\lor\;\; \{u_a,v_a\}\cap\mathcal{P}\neq\emptyset,
\label{eq:builder}
\end{equation}
where $\mathcal{P}$ is the set of verified protocol contracts; on acceptance $d(v_a)\leftarrow\min(d(v_a),d(u_a)+1)$ unless $v_a\in\mathcal{P}$, so the expansion never walks \emph{through} shared infrastructure.
The collector then fetches the new frontier $\{v_a\}\setminus\mathcal{P}$ and the builder is re-run until the frontier is empty.
We use $D=6$ hops, $H=72$\,h and $\eta=5\%$ for incidents and $D=2$ for negatives, whose seeds are ordinary users.
The value share is computed per token so that a large stablecoin payment is not diluted by a small ETH fee.
Traces are followed on the incident's primary chain; the continuation on the destination chain of a bridge deposit is not stitched into the trajectory, because the reports we rely on confirm the exit but seldom give the destination address.
This is why a bridge deposit is a natural endpoint in our formulation.

\subsection{Prefix-Safe Features}
\label{sec:features}
For each prefix the extractor computes $\mathbf{x}(\tau_{1:k})$ from $a_1,\dots,a_k$ only (37 model inputs in total).
A unit test enforces this: two trajectories that share the first $k$ actions but differ afterwards, in content or in length, must yield identical features at $k$.
The same test caught a \emph{relative-position} feature $k/n$ in an earlier version of the code, whose denominator is the final length; it is removed, and every result in this paper is computed without it.

\begin{itemize}
\item \textbf{Flat (14).} Temporal: mean and standard deviation of inter-action gaps $\delta_i=t_{i}-t_{i-1}$, burstiness $B=(\sigma_\delta-\mu_\delta)/(\sigma_\delta+\mu_\delta)$~\cite{goh2008burstiness}, active duration $t_k-t_1$.
Structural: fan-out and fan-in of the seed, number of distinct counterparties, edge density $|E|/(|V|(|V|-1))$, maximum hop depth from the seed, number of branching addresses (out-degree $>1$).
Economic: mean of $\log(1+q_i)$, out/in value ratio, value retention $q_k/q_1$, and the number of tokens holding at least 1\% of the observed value.
\item \textbf{Typed actions (16).} The share $\frac1k|\{i\le k:y_i=y\}|$ of each type $y\in\mathcal{Y}$, the number of distinct type bigrams and trigrams, the number of distinct bridge families, the number of bridge$\to$swap transitions, the number of stablecoin legs, and the time to the first bridge action.
Raw counts are not used because they sum to $k$ and would re-introduce prefix length as a shortcut.
\item \textbf{Motifs (7).} Counted only when complete within the prefix. \emph{Split}: addresses with $\ge2$ outgoing actions; \emph{merge}: addresses with $\ge2$ incoming actions; \emph{peel chain}: a run of $\ge3$ consecutive transfers $a_{i-1}\to a_i$ with $v_{i-1}=u_i$ and $q_i\le q_{i-1}$; \emph{bridge$\to$swap} and \emph{swap$\to$split}: the ordered pair in consecutive actions; \emph{nested bridge}: $\ge2$ distinct bridge protocols; \emph{rapid token pivot}: the number of token changes between consecutive actions.
\end{itemize}
The flat group is exactly the information a type-agnostic transaction graph offers; the other two groups require the decoder.
Table~\ref{tab:aadapt} maps the observables to MITRE AADAPT~\cite{mitre2025aadapt} techniques so that an alert can be phrased in a responder's vocabulary.

\begin{table}[t]
\centering
\caption{CAGI-ED observables mapped to MITRE AADAPT techniques.}
\label{tab:aadapt}
\setlength{\tabcolsep}{4pt}
\small
\begin{tabular}{@{}llp{5.1cm}@{}}
\toprule
AADAPT ID & Technique & CAGI-ED observable \\
\midrule
ADT3005 & Cross-Chain Swaps (Hopping) & bridge deposit share, bridge$\to$swap, nested bridge \\
ADT3028 & Siphon Funds & split, fan-out, distinct counterparties \\
ADT3028.001 & Anonymous Wallets & hop depth through fresh intermediaries \\
ADT3028.003 & Layering & swap share, rapid token pivot, bigram diversity \\
ADT3028.005 & Peel Chains & peel chain, value retention \\
ADT3030 & Use Anonymizing Services & mixer/exit action (endpoint) \\
\bottomrule
\end{tabular}
\end{table}

\subsection{Scoring, Calibration and Explained Alerts}
\label{sec:scoring}
The scorer $f$ is a gradient-boosted tree ensemble (XGBoost~\cite{chen2016xgboost}, 200 trees of depth 4, learning rate 0.1, class weight $n_-/n_+$ computed on the training fold), fixed before any experiment and never tuned on the evaluation data.
Inside each training fold, features present in fewer than two positive incidents are dropped, so a feature cannot memorize one incident.
Raw scores are mapped to probabilities by Platt scaling~\cite{niculescu2005calibration},
\begin{equation}
\hat p_k=\sigma\!\bigl(\alpha\,\mathrm{logit}\,f(\mathbf{x}(\tau_{1:k}))+\beta\bigr),
\label{eq:platt}
\end{equation}
with $(\alpha,\beta)$ fitted on out-of-fold scores of the training incidents only.
When $\hat p_k\ge\theta$ for the first time, the alert carries the $k$ with the largest exact TreeSHAP contributions~\cite{lundberg2020treeshap} $\phi_j(\mathbf{x})$, where $\mathrm{logit}\,f(\mathbf{x})=\phi_0+\sum_j\phi_j(\mathbf{x})$, together with the AADAPT IDs of their features.
Because features are recomputed from the prefix, scoring is stateless and can be repeated on every block.
